\paragraph{Qué reproduce el modelo.} Con la ciudad completa, la configuración seleccionada para la prueba reservada (z02) reproduce los dos corredores que delimitan la zona afectada oficial ---Camino El Alba y General Blanche por el norte, Quebrada Honda por el sur---, el alcance hasta la Av.\ Padre Hurtado, cerca de un metro de barro en la urgencia de la Clínica UC Christus y la afectación de Duoc UC y de la Clínica Universidad de los Andes. El volumen de entrada que reproduce ese alcance en el modelo, unos \SI{300000}{m^3}, coincide con el que se obtiene de forma independiente a partir de la lluvia informada, suponiendo suelo saturado y lluvia en toda la superficie de la cuenca. La concordancia entre estimaciones obtenidas mediante procedimientos distintos aporta coherencia al orden de magnitud, bajo sus respectivos supuestos; no constituye validación independiente.

\paragraph{Lluvia antecedente, isoterma y volumen.} Según T13, el 31 de enero cayeron \SI{75}{mm} en \SI{80}{min}, y su balance de afectación fue menor que los informados para octubre. Las intensidades informadas para octubre son menores; la alcaldesa informó unos \SI{144}{mm} acumulados y el CR2 situó la isoterma de \SI{0}{\celsius} por sobre la cumbre de la cuenca. Estos antecedentes motivan escenarios con distintos supuestos de saturación y área bajo lluvia; la comparación no establece la causa del evento ni la contribución relativa de esos factores. Bajo los supuestos del cálculo, el límite superior del aporte por deshielo es menor al \SI{5}{\percent}: no se detectó nieve en la parte despejada de la cuenca en la imagen Sentinel-2 del 21 de septiembre.

\paragraph{El papel del drenaje (escenarios hipotéticos).} En los escenarios ensayados, la afectación de Quebrada Honda se reproduce cuando parte del flujo calculado no ingresa al colector de Francisco Bulnes Correa. Con la ciudad completa, el caudal calculado en la embocadura llega a unos \SI{65}{m^3/s}; bajo los supuestos ensayados, una capacidad supuesta de \SI{30}{m^3/s} deja flujo superficial y el análisis sitúa cerca de \SI{65}{m^3/s} el umbral calculado de capacidad supuesta para retirarlo por completo, con captura de toda la sección en el modelo. Estos valores no describen la obra real. La cartografía hidrográfica oficial representa la continuidad del cauce por Quebrada Honda aguas abajo de la embocadura; en las simulaciones, el flujo que no ingresa al colector continúa por ese trazado. El análisis delimita capacidades supuestas bajo los supuestos del modelo; no determina la capacidad, el estado ni el funcionamiento de la obra real durante el evento y no permite atribuir responsabilidades.

\paragraph{Posibles explicaciones de la extensión simulada fuera del perímetro oficial.} En z02, solo un tercio del barro urbano simulado cae dentro de la zona afectada oficial; con casas porosas (z26), la fracción sube a 0,54, y con la hipótesis integrada corregida y lluvia medida (p34b), a 0,48. La referencia es la zona afectada oficial, que es un perímetro de gestión y no una huella verificada del barro. Se ensayaron las siguientes representaciones; sus resultados describen sensibilidad del modelo y su correspondencia con el evento permanece sin probar:
\begin{enumerate}
\item \textbf{Edificios impermeables.} En la configuración inicial cada edificio es un bloque macizo; las fuentes de prensa recogidas en el relato informan ingreso de barro a viviendas. Se ensayó como supuesto alternativo la entrada y retención en los lotes mediante edificios porosos, una representación habitual en inundación urbana \citep{schubert2012} (v26, verificada con paso fino en z26). La fracción dentro de la zona afectada oficial sube de 0,34 a 0,54, pero el espesor junto a la Clínica UC Christus baja a \SI{0,30}{m}.
\item \textbf{Carga gruesa.} Las imágenes consultadas, descritas en el marco teórico, muestran piedras, troncos y vehículos arrastrados. La reología de la matriz fina no incluye esa resistencia adicional; un barro algo más concentrado ($C_v=0{,}43$, entre el que se extiende y el que se queda corto) la representó de forma simple (v25), sin reducir la extensión fuera de la zona afectada oficial (fracción dentro: 0,346 frente a 0,344 en z02).
\item \textbf{Vegetación y jardines.} FABDEM elimina árboles y la rugosidad fuera de las calles es moderada; al aumentarla (v27), la fracción dentro de la zona afectada oficial sube a 0,388.
\end{enumerate}
A esto se suma una limitación de la métrica: el modelo no distingue el agua de la lluvia del barro del aluvión, de modo que la escorrentía de la lluvia en otros barrios cuenta como barro fuera de la zona afectada oficial y penaliza a las hipótesis con lluvia en la ladera.

\paragraph{Sensibilidad al terreno.} En los ensayos, la representación explícita de los cauces menores es lo que más cambia el recorrido calculado: el tallado del cauce del parque lineal modifica la presencia de flujo simulado en Quebrada Honda, el trazado supuesto del brazo sur modifica su llegada a la parte alta de Álvaro del Portillo y la forma de representar la embocadura cambió en un \SI{44}{\percent} la cobertura de Quebrada Honda. Con un terreno de \SI{30}{m}, esos corredores deben tallarse a partir de la cartografía; la red hidrográfica oficial a escala 1:25.000 resultó más fiable que OpenStreetMap para el brazo sur.

\paragraph{Campus y patio hundido.} El desborde simulado en Álvaro del Portillo desciende hacia General Blanche y no llega al patio con ninguno de los fluidos ensayados. Las variantes de escorrentía local, fluido más líquido y representación del claro y de la calzada baja de La Plaza Sur sí se completaron: las microsimulaciones y m36 reproducen la entrada La Plaza Sur $\rightarrow$ claro $\rightarrow$ patio. Con malla fina, m36 calcula \SI{0,17}{m} en el patio y moja el \SI{58}{\percent} de La Plaza Sur (sección~\ref{sec:res-posterior}). Este resultado es compatible con las imágenes consultadas, pero no identifica el proceso ocurrido durante el evento ni constituye validación independiente.

\paragraph{Prueba con observaciones reservadas.} La única prueba con observaciones reservadas, definida antes de cruzarlas con el modelo, fue desfavorable: la hipótesis seleccionada alcanza 19 de 49 equipamientos afectados reservados (\SI{39}{\percent}), frente al \SI{70}{\percent} fijado, y no supera a la hipótesis de solo el cauce principal (sección~\ref{sec:res-ronda}). El IEC de 72/100 mide compatibilidad con observaciones usadas para formular las hipótesis; esta prueba muestra que esa compatibilidad no se traslada a la distribución espacial fina de la afectación oficial. El modelo también deja sin barro partes importantes de la zona afectada oficial, no solo la excede. Esto puede deberse a la resolución del terreno, a la representación de los edificios, al volumen o su reparto en el tiempo, a aportes locales no modelados, o a diferencias entre la definición de afectación de la capa oficial y el umbral de espesor modelado. La representación de los edificios, la resolución local, el volumen y la distribución temporal de la lluvia se ensayaron después (sección~\ref{sec:res-posterior}); esos ajustes no cambian el resultado de la prueba reservada y su validación independiente permanece sin probar.

\paragraph{Lo que mejoró después de la prueba.} Dos cambios de representación explican la mayor parte de la mejora posterior a la prueba reservada. El primero, permitir que el barro entre a los lotes, redistribuye la mezcla dentro de la zona afectada oficial. El segundo consiste en suponer agua turbia en el campus y La Plaza Sur y barro espeso en la rama principal, y en representar los caminos y las calzadas encajonadas que conducen el agua. Ambos tienen sustento observacional (videos, Street View, Google Earth). Aun así, se eligieron conociendo los resultados, de modo que el aumento del IEC (de 72 a 76--82 en los ajustes verificados con paso fino) mide un mejor ajuste y no una validación. Al cierre, la hipótesis integrada obtiene 79--82: 82 con \SI{35}{mm/h} supuestos y paso fino (z32), 81 con el pulso medido en Río Molina (p34b, corregida) y 79 con esa lluvia repartida en \SI{4}{h} (v33). La disminución del espesor en la Clínica UC Christus muestra que la retención en los lotes, representada así, sobrestima la pérdida de espesor en la calle.

\paragraph{Incendio y métricas.} La severidad estimada con dNBR (23 de diciembre de 2025 y 6 de febrero de 2026) da un \SI{82}{\percent} con señal de quema y un \SI{10}{\percent} de severidad moderada-alta o alta; con esta última fracción, la mediana de Gartner es de \SI{148000}{m^3}. Con los coeficientes supuestos, el balance sin incendio calcula alrededor del \SI{62}{\percent} del volumen y el \SI{43}{\percent} de los sólidos; Staley M1 da 0,28 con incendio frente a 0,07 sin él, como uso ilustrativo. Las repeticiones corregidas f35b y f38b mojan \SIrange{179}{193}{ha}, frente a \SI{222}{ha} en p34b (\SIrange{13}{19}{\percent} menos), y calculan \SIrange{0,19}{0,23}{m} junto a Clínica UC Christus. El IEC de 80--82, frente a 81, muestra que este índice no distingue esos escenarios (sección~\ref{sec:res-incendio}).

\paragraph{Validez.} Las observaciones disponibles son cualitativas: calles y recintos informados por la prensa y las autoridades, y un perímetro oficial que delimita la zona de gestión y no el área con barro verificado. No hay aforos ni marcas de nivel publicadas. El índice exploratorio de compatibilidad resume cuatro criterios fijados antes de simular y sirve para comparar hipótesis entre sí; no es una medida de precisión predictiva, y diferencias de dos o tres puntos no deben interpretarse como evidencia concluyente de superioridad entre hipótesis.

\paragraph{Grado de confianza por resultado.} En lugar de reducir los resultados a un único indicador, la tabla~\ref{tab:confianza} asigna a cada uno una valoración cualitativa de confianza al cierre.
\begin{table}[htbp]\centering\small
\caption{Valoración cualitativa de la confianza en cada resultado al cierre, condicionada por los supuestos y la evidencia disponible.}\label{tab:confianza}
\begin{tabular}{p{0.42\textwidth}p{0.5\textwidth}}\toprule
Resultado & Grado de confianza\\\midrule
Corredores generales del barro & Relativamente alta: compatibles con múltiples escenarios y con la zona afectada oficial\\
Necesidad de representar cauces y transiciones hacia colectores & Hallazgo sólido de sensibilidad del modelo\\
Orden de magnitud del volumen movilizado & Estimación condicionada, con rango amplio (\SI{181000}{m^3} a \SI{468000}{m^3})\\
Capacidad real del colector & No determinada; solo capacidades ensayadas en escenarios hipotéticos\\
Profundidad exacta frente a edificios & Exploratoria, dependiente del terreno y de los parámetros\\
Velocidades, presiones y tiempos de llegada & Indicadores calculados, sin validación instrumental\\
\bottomrule\end{tabular}
\end{table}
